% Generated draft: requires mathematical and source review.
\documentclass[11pt]{article}
\usepackage[margin=25mm]{geometry}
\usepackage{fontspec}
\setmainfont{Latin Modern Roman}
\usepackage{amsmath,amssymb,mathtools,unicode-math}
\setmathfont{Latin Modern Math}
\usepackage{enumitem,needspace,xurl,hyperref}
\setlength{\parindent}{0pt}
\setlength{\parskip}{5pt}
\setlength{\emergencystretch}{3em}
\newcommand{\N}{\mathbb{N}}
\newcommand{\Z}{\mathbb{Z}}
\newcommand{\Q}{\mathbb{Q}}
\newcommand{\R}{\mathbb{R}}
\newcommand{\C}{\mathbb{C}}
\newcommand{\F}{\mathbb{F}}
\newcommand{\A}{\mathbb{A}}
\newcommand{\PP}{\mathbb{P}}
\newcommand{\E}{\mathbb{E}}
\newcommand{\eps}{\varepsilon}
\newcommand{\dd}{\,\mathrm{d}}
\begin{document}
\section*{UnsolvedMath ID: 800920}
\textbf{Source problem code:} catalogue entry
\subsection{Definitions and mathematical statement}
Let \(S\) be a \(K3\) surface, let \(g\in\Z_{\geq 0}\) denote the genus, and let \(\beta\) denote a curve class on \(S\). The ground field, the group containing \(\beta\), and the admissibility conditions on curve classes are not specified in the supplied record.

\textbf{Requested assertion (Pandharipande, 2010, Q8).} Prove the Katz--Klemm--Vafa conjecture in every genus and in every class on \(K3\) surfaces, with the conventions and hypotheses intended in the cited source.

The supplied record names the conjecture but does not give its identity, define the invariants occurring in it, or specify their normalization. Thus this is a named-conjecture statement only: a self-contained mathematical formula cannot be recovered from the supplied text. In particular, the phrase ``all classes'' is retained without silently restricting it to primitive, effective, or nonzero classes; its precise scope requires clarification from the source.

The accompanying source commentary presents a solution as providing a large class of exact formulas for higher-genus Gromov--Witten invariants of compact Calabi--Yau threefolds. It also describes mathematical results in that setting as difficult to obtain, in contrast with local toric cases, and refers to an unresolved reference label \textup{[zing]} for the genus-one theory of the quintic threefold. These are supplied historical motivations, not additional assertions of the conjecture or an independently verified account of the literature.
\subsection{Short English statement}
Pandharipande's 2010 question asks for the Katz--Klemm--Vafa conjecture for \(K3\) surfaces in every genus and every curve class. The supplied excerpt omits the conjectural formula and essential conventions, so this draft preserves the named target and explicitly records the missing formulation. The registry labels the problem open; that label has not been independently verified.
\subsection{Sources}
\begin{itemize}
\item Local UnsolvedMath registry: permanent ID \(800920\), exact source problem code \texttt{catalogue entry}, title ``AGNES Open Problems in Algebraic Geometry --- Q8.'' The supplied record attributes the question to Pandharipande, with proposed year 2010, and labels its status \texttt{open}.
\item Supplied source attribution: Pandharipande, \emph{Open Problems in Enumerative Algebraic Geometry (AGNES)} (2010), source item Q8 (\#7), \url{https://arxiv.org/abs/1004.3259}. The provenance records access on 2026-10-16 and extraction method \texttt{arxiv\_latex}; these are supplied provenance claims, not independent source verification.
\item The source excerpt cites \textup{[zing]} for the genus-one theory of the quintic threefold. No corresponding bibliographic entry or URL is supplied, so this reference cannot be identified reliably from the record.
\end{itemize}
\end{document}
